problem:  
Let \((M^n,g(t))_{t\in(-\infty,0]}\) be a complete κ-noncollapsed ancient solution to the Ricci flow with bounded, nonnegative curvature operator, and fix a basepoint \(p\in M\).  
For each \(t<0\), let \(l_t(x):=l(x,t)\) be Perelman’s reduced distance based at \((p,0)\), and define the **reduced distance minimizer set**
\[
C_t := \{ x \in M : l_t(x) = \min_M l_t \}.
\]
By standard growth estimates on \(l_t\), each \(C_t\) is nonempty and contained in a compact region.

Suppose that for some sequence \(t_k\to -\infty\), the rescaled flows \((M,|t_k|^{-1}g(|t_k|s),p)\) subconverge in Cheeger–Gromov sense to a nonflat gradient shrinking soliton \((M_\infty,g_\infty(s),f_\infty,p_\infty)\).  
Let \(C_\infty\subset M_\infty\) denote the set where the soliton potential \(f_\infty\) achieves its minimum.

**Open problem:**  

1. Does there exist a universal geometric mechanism ensuring that, after rescaling, the sets \(C_{t_k}\) converge (e.g., in Hausdorff distance under pullback by convergence diffeomorphisms) to \(C_\infty\)?  

2. More profoundly, is this limiting minimizer set **independent of the chosen backward sequence** \(t_k\to-\infty\)?  
Equivalently, does every κ-noncollapsed ancient solution possess a well-defined *asymptotic reduced-distance center*—a canonical geometric locus attracting all reduced-distance minimizers under backward evolution?  

If not, can one expect ancient flows that admit more than one distinct family of minimizing loci corresponding to different asymptotic gradient shrinking solitons or their translations?

---

Why is it a "good" problem:  

This version is both **intrinsically formulated** and **conceptually rich**, and avoids coordinate dependency while targeting a subtle, open rigidity question central to the backward-time geometry of Ricci flow.

1. **Intrinsic and mathematically clean:**  
   The reduced-distance function \(l(x,t)\) and its minimizers are canonical objects in Perelman’s entropy framework. The problem uses them to track geometric centers without invoking external coordinates or embeddings.

2. **Beyond uniqueness of tangent flows:**  
   While related to questions of blow-down uniqueness, this problem isolates a **geometric localization phenomenon**—it asks whether the *points* (or sets) where backward time is least costly to reach converge to a unique “center,” independent of rescaling sequences. Proving or disproving such stabilization represents a new kind of rigidity, complementing metric-soliton uniqueness results.

3. **Analytic–geometric bridge:**  
   The question interfaces Perelman’s reduced geometry (analysis of \(l\)-functions, differential Harnack inequalities) with limit-flow theory (Cheeger–Gromov convergence of solitons). Understanding convergence of \(C_t\) could require new insights into parabolic potential theory and the fine structure of the reduced distance under rescaling.

4. **Potential implications:**  
   - A positive answer would establish existence of an **intrinsic backward center** for any ancient flow, suggesting new canonical normalizations for blow-down analysis and potentially strengthening uniqueness results for asymptotic solitons.  
   - A counterexample, with multiple limit centers, would exhibit previously unseen **drift behavior** of entropy concentration, implying that ancient flows may have multi-centered or dynamical asymptotic structures.

5. **Simplicity of statement, depth of content:**  
   The question is succinct—“Do reduced-distance minimizers stabilize or drift as \(t\to -\infty\)?”—but its resolution touches on profound analytic and geometric mechanisms underlying Ricci flow rigidity and the emergence of canonical self-similar models.

Hence, it presents an elegant, intrinsic, and open challenge: determining whether the backward geometric core of an ancient flow possesses a unique and stable center of shrinkage. This blend of simplicity, natural formulation, and potential theoretical reach makes it a “good” research-level problem.